\documentclass[a4paper,12pt]{article}
\usepackage[margin=2.1cm,headheight=16pt]{geometry}
\usepackage{xcolor,booktabs,array,graphicx,fancyhdr}
\usepackage{hyperref}
\hypersetup{colorlinks=true,urlcolor=teal,linkcolor=teal,pdftitle={مقایسه پژوهش‌های مرتبط با BlunderGuard},pdfauthor={BlunderGuard}}
\usepackage{xepersian}
\settextfont[AutoFakeSlant=0.15]{Noto Sans Arabic}
\setlatintextfont{DejaVu Sans}
\setdigitfont{Noto Sans Arabic}
\linespread{1.18}
\setlength{\emergencystretch}{5em}
\renewcommand{\labelitemi}{\lr{\textbullet}}
\setlength{\parindent}{0pt}
\setlength{\parskip}{6pt}
\pagestyle{fancy}
\fancyhf{}
\fancyhead[R]{مقایسهٔ پژوهش‌های مرتبط}
\fancyhead[L]{\lr{BlunderGuard}}
\fancyfoot[C]{\thepage}
\newcommand{\en}[1]{\lr{#1}}
\newcommand{\pathref}[1]{\begin{latin}\small\texttt{\detokenize{#1}}\end{latin}}
\newcommand{\source}[2]{\par\small منبع اصلی: \href{#1}{#2}\par\normalsize}
\newcommand{\paper}[3]{\clearpage\section{#1}\begin{latin}\small #2\end{latin}\textbf{وضعیت انتشار:} #3\par}
\begin{document}
\begin{titlepage}
\vspace*{2cm}
\begin{center}
{\Huge\bfseries پروژهٔ من در برابر پژوهش‌های مرتبط}\\[1cm]
{\LARGE مقایسهٔ تحلیلی \en{BlunderGuard}}\\[0.5cm]
{\Large با هشت مقالهٔ پژوهشی و یک مقالهٔ فنی}\\[1.5cm]
{\large تولید توضیح شطرنج، آموزش مدل کوچک، اتصال به موتور و سنجش صحت}\\[1.5cm]
\end{center}
\textbf{پرسش اصلی سند:} پروژه در مقایسه با هر کار چه قابلیت اضافه‌ای دارد، چه چیزی کم دارد، و برای اثبات برتری خود به چه آزمایشی نیاز دارد؟

\textbf{مبنای بررسی:} کد، نوت‌بوک آموزشی، گزارش و نتایج موجود در مخزن، به‌علاوه متن منابع اصلی. تاریخ بررسی: \en{2026-09-13}.

\textbf{نتیجهٔ کلیدی:} مزیت قابل‌دفاع فعلی، ترکیب یک مولد کوچک محلی با شواهد صریح موتور و رابط کاربردی است. برتری در صحت توضیح، کاهش توهم یا اثر آموزشی هنوز اثبات نشده است؛ نتایج داخلی فعلی نیز ضعف مدل اختصاصی در انطباق با مرجع را نشان می‌دهند.

\vfill
این سند گزارش مقایسه و تحلیل است؛ نتیجهٔ اجرای مجدد آموزش، بازتولید کامل مقاله‌ها یا آزمایش مستقیم میان سامانه‌ها نیست.
\end{titlepage}
\tableofcontents
\clearpage
\section{روش مقایسه و حدود اعتبار}
این بررسی هر ۹ منبع معرفی‌شده در جست‌وجوی قبلی را پوشش می‌دهد. ترتیب منابع بر اساس نزدیکی به هدف پروژه است، نه رتبه‌بندی کیفیت مقاله‌ها. خلاصهٔ روش‌ها از متن اصلی خوانده شده و قضاوت دربارهٔ پروژه بر شواهد محلی تکیه دارد.

\subsection*{چهار نوع گزاره را از هم جدا می‌کنیم}
\begin{itemize}
\item \textbf{قابلیت موجود:} در کد یا خروجی ذخیره‌شده شاهد دارد؛ وجود کد الزاماً به معنی کارکرد بی‌نقص نیست.
\item \textbf{مزیت طراحی:} انتخابی که برای هدف مشخصی مفید است، مثل استقلال از سرویس ابری؛ به‌تنهایی برتری علمی یا سرعت بیشتر را ثابت نمی‌کند.
\item \textbf{کمبود شواهد:} آزمایش، داده یا قابلیت لازم در بخش‌های بررسی‌شده پیدا نشده است؛ این با اثبات نبود آن در همه نسخه‌های پروژه یا مقاله فرق دارد.
\item \textbf{برتری اثبات‌شده:} فقط با اجرای مقایسهٔ منصفانه روی داده و معیار مشترک قابل ادعاست. چنین مقایسه‌ای با سامانه‌های این مقاله‌ها در مخزن بررسی‌شده یافت نشد.
\end{itemize}

\subsection*{محدودیت‌های ضروری برای تفسیر}
امتیازهای مقاله‌های مختلف مقیاس و مجموعه‌آزمون یکسان ندارند. امتیاز داور زبانی پروژه را نمی‌توان با دقت حل مسئله یا امتیاز انسانی یک مقاله مستقیماً مقایسه کرد. جدیدتر بودن مدل، استفاده از چند مسیر موتور یا داشتن رابط وب نیز به‌خودی‌خود نوآوری پژوهشی محسوب نمی‌شود.

عبارت «در مقاله گزارش نشده» در این سند به معنی «آن سامانه قطعاً ندارد» نیست. به‌ویژه برای اجرای محلی و رابط کاربری، تمرکز مقاله ممکن است بر بخش دیگری باشد. نسخهٔ \en{CARLSy} ناشناس است و پذیرش آن تأیید نشده؛ مقالهٔ \en{Tredence} نیز وبلاگ فنی است، نه پژوهش داوری‌شده. سایر پیش‌چاپ‌ها با همان وضعیت آرشیوی معرفی می‌شوند.

برای تدوین، کد تغییر نکرده و مدل زبانی یا موتور در یک ارزیابی کامل اجرا نشده است. تنها یک بررسی مستقل و سبک برای خطای قالب پرامپت انجام شد. توصیه‌های این سند پیشنهاد پژوهشی‌اند و به‌عنوان قابلیت فعلی شمرده نمی‌شوند.

\clearpage
\section{پروژه واقعاً چه دارد؟}
\subsection*{معماری و آموزش}
ورودی اصلی دو وضعیت \en{FEN} قبل و بعد از حرکت است. کد موتور برای هر وضعیت سه ادامهٔ برتر را در عمق ۱۷ می‌گیرد. استخراج‌گر صفحه، مهرهٔ تحت حملهٔ بدون مدافع، تغییر امنیت شاه، کنترل مرکز، گسترش و برخی ضعف‌های پیاده را توصیف می‌کند. این قواعد قابل‌خواندن هستند؛ اما صرف قطعی بودن اجرای آن‌ها به معنی کامل یا درست بودن تعریف شطرنجی آن‌ها نیست.

نوت‌بوک، مدل پایهٔ \en{Qwen3-1.7B}، سازگارکنندهٔ \en{LoRA} با رتبهٔ ۱۶، آموزش زبان علّی روی قطعات ۲۵۶ توکنی متن کتاب و خروجی \en{GGUF Q8\_0} را نشان می‌دهد. این آموزش به‌طور مستقیم نگاشت «وضعیت و حرکت به توضیح درست» را هدف نمی‌گیرد. مقدار بیش از ۲٫۸ میلیون توکن در گزارش آمده و در این بررسی مستقلاً بازشماری نشده است.

\subsection*{تولید و محصول}
کد دو روش دارد: یک فراخوانی مستقیم و روش زنجیره‌ای با دو توضیح میانی و یک جمع‌بندی. پرامپت نهایی توضیح سادهٔ دو یا سه جمله‌ای می‌خواهد و نام حرکت، نمادگذاری و عدد را منع می‌کند. مسیر سرویس محلی \en{llama.cpp}، بک‌اند \en{Flask} و فرانت‌اند \en{React} موجود است؛ ابزار تولید دسته‌ای نیز مدل محلی یا سرویس سازگار با \en{OpenAI API} می‌پذیرد.

\subsection*{ارزیابی}
خلاصهٔ منتخب شامل هفت مدل و برای هر مدل ۶۰ نمونه در سه گروه ۲۰تایی است. طبق گزارش، گروه‌بندی از سطح بازیکنان بازی‌ها می‌آید؛ این روش به‌تنهایی دشواری ذاتی هر موقعیت یا شخصی‌سازی متن برای کاربر را ثابت نمی‌کند. داور، پاسخ و متن مرجع را می‌گیرد و پنج بُعد را می‌سنجد. در این مسیر، داور به صفحه و ابزار موتور برای بررسی ادعا دسترسی ندارد.

\textbf{شواهد محلی اصلی:}
\pathref{app/chess/stockfish_analyzer.py | app/chess/position_analysis.py}
\pathref{app/chat.py | app/prompt.py | app/server.py}
\pathref{notebooks/chess.ipynb | evaluation/judge_responses.py}
\pathref{docs/report/main.tex | evaluation/judge_gemma12b/summary.csv}

\clearpage
\section{نتایج داخلی؛ مزیت و ضعف را با هم ببینیم}
اعداد زیر میانگین وزنی سطرهای فایل خلاصهٔ منتخب هستند؛ وزن هر گروه تعداد نمونهٔ آن است. در این فایل همهٔ گروه‌ها ۲۰ نمونه دارند و خطای تجزیهٔ ثبت‌شده صفر است. نام مدل‌ها عین شناسهٔ خروجی‌هاست و هویت سرویس‌های میزبان مستقلاً راستی‌آزمایی نشده است.

\begin{center}\small
\begin{tabular}{p{6.7cm}rrr}
\toprule
مدل & صحت & روانی & کل \\
\midrule
\en{chess-q8\_0} & ۰٫۸۳ & ۸٫۹۷ & ۲٫۴۹ \\
\en{Llama-3.2-3B-Instruct-Q4\_K\_M} & ۰٫۸۳ & ۱۰٫۰۰ & ۲٫۸۶ \\
\en{gemma-3-12b-it-Q4\_0} & ۰٫۸۸ & ۱۰٫۰۰ & ۲٫۹۵ \\
\en{grok-4.3} & ۳٫۹۲ & ۱۰٫۰۰ & ۵٫۴۴ \\
\en{gemini-3.5-flash} & ۵٫۱۵ & ۱۰٫۰۰ & ۶٫۴۶ \\
\en{qwen3.7-plus} & ۵٫۳۰ & ۱۰٫۰۰ & ۶٫۵۰ \\
\en{gpt-5.5} & ۶٫۲۸ & ۱۰٫۰۰ & ۷٫۳۱ \\
\bottomrule
\end{tabular}
\end{center}
همهٔ امتیازها از ۱۰ هستند. وزن امتیاز کل: صحت ۳۵ درصد، وفاداری ۲۵ درصد، سودمندی ۲۰ درصد، پیروی از دستور ۱۰ درصد و روانی ۱۰ درصد.

\textbf{برداشت قابل‌دفاع:} مدل اختصاصی در این پروتکل روان است، اما انطباق معنایی و واقعی ضعیفی با مرجع دارد. دادهٔ فعلی ادعای «آموزش روی کتاب، مدل کوچک را از مدل‌های عمومی بهتر کرده» را پشتیبانی نمی‌کند.

\textbf{برداشت غیرمجاز:} صحت ۰٫۸۳ از ۱۰ معادل ۸٫۳ درصد ادعای درست نیست. داور صریحاً مرجع را تنها حقیقت مبنا می‌گیرد و نباید صفحه را استنتاج کند؛ پس امتیاز، توافق با مرجع است و می‌تواند توضیح درستِ متفاوت با مرجع را جریمه کند. این عدد نرخ واقعی توهم نیز نیست.

\textbf{کمبود کنترلی:} بدون نتیجهٔ مدل پایهٔ هم‌اندازه با همان پرامپت و کمّی‌سازی نمی‌توان اثر آموزش را جدا کرد. صحت‌سنجی انسانی، بازهٔ اطمینان و مقایسهٔ مستقیم با مقاله‌ها نیز در این خلاصه وجود ندارد. خروجی داور دیگر در مخزن موجود است، اما به دلیل تفاوت اجرا و خطاهای ثبت‌شده در این جدول ادغام نشده است.
\pathref{evaluation/judge_gemma12b/summary.csv}

\paper{مقایسه با \en{Kim et al.}}{Bridging the Gap between Expert and Language Models: Concept-guided Chess Commentary Generation and Evaluation}{\en{NAACL 2025}؛ پیش‌چاپ اولیه ۲۰۲۴}
\textbf{روش مقاله.} چارچوب \en{CCC} با برچسب‌های مفهومی موتور کلاسیک، جهت‌های مفهومی را از بازنمایی \en{LeelaChessZero} استخراج می‌کند؛ تغییر اهمیت مفاهیم قبل و بعد از حرکت، راهنمای مولد است. \en{GCC-Eval} ارزیابی چندبعدی را با شواهد تخصصی انجام می‌دهد و با قضاوت انسانی سنجیده شده است. بخش‌های ۳ و ۴ مبنای این مقایسه‌اند.

\textbf{وجه مشترک.} هر دو روش می‌خواهند شواهد تخصصی شطرنج را به زبان قابل‌فهم تبدیل کنند و تغییر موقعیت را برای انتخاب محتوای توضیح به کار بگیرند. بنابراین ایدهٔ «مفهوم شطرنجی به‌عنوان واسط موتور و زبان» مختص پروژه نیست.

\textbf{پروژه چه چیز اضافه یا متفاوتی دارد؟} در مسیر فعلی، مولد کوچک کمّی‌شده روی دستگاه محلی اجرا می‌شود. استخراج ویژگی‌های دستی به آموزش طبقه‌بند روی بازنمایی داخلی موتور نیاز ندارد و بررسی هر قاعده آسان‌تر است. این یک مزیت سادگی و استقلال است؛ در برابر استفاده از مدل‌های اختصاصی بزرگ در آزمایش اصلی مقاله، مسیر محلی پروژه تفاوت عملی مشخصی دارد.

\textbf{چه چیز کمتر دارد؟} پروژه اولویت‌بندی آموخته‌شدهٔ مفاهیم یا سنجش ارتباط آن‌ها با بازنمایی تصمیم موتور ندارد. همچنین تأیید انسانی معتبر برای داور و آزمایش مستقیم با \en{CCC} ارائه نشده است. پنج گروه ویژگی دستی پروژه جای پوشش گسترده و انتخاب سنجیدهٔ محتوای توضیح را تضمین نمی‌کنند.

\textbf{آزمایش تعیین‌کننده.} با یک مولد ثابت، سه ورودی «فقط موتور»، «موتور و قواعد پروژه» و «مفاهیم اولویت‌بندی‌شده» را روی موقعیت‌های یکسان بسنجید. سپس مولد کوچک را با مولد بزرگ در هزینه و صحت مقایسه کنید. در وضعیت فعلی مزیت استقرار قابل‌طرح است، ولی برتری توضیح ثابت نشده است.
\source{https://aclanthology.org/2025.naacl-long.481/}{صفحهٔ انتشار و متن کامل مقاله}

\paper{مقایسه با \en{Lee et al.}}{Improving Chess Commentaries by Combining Language Models with Symbolic Reasoning Engines}{پیش‌چاپ \en{arXiv}، دسامبر ۲۰۲۲}
\textbf{روش مقاله.} مدل زبانی کنترل‌پذیر با برچسب‌هایی دربارهٔ نوع توضیح، کیفیت حرکت، حرکت پیشنهادی، ارجاع‌های زبانی و طول متن آموزش می‌بیند. موتور نمادین در استنتاج اطلاعات کنترلی را تأمین می‌کند. روش بر دادهٔ جفت‌شدهٔ وضعیت و توضیح تکیه دارد و ترجیح داور انسانی را گزارش می‌کند؛ بخش ۴ روش و بخش آزمایش‌ها مبنا هستند.

\textbf{وجه مشترک.} این نزدیک‌ترین پیشینهٔ معماری است: محاسبهٔ شطرنج به موتور سپرده می‌شود و بیان آن به مدل زبان. محدود کردن خروجی و هدایت کیفیت حرکت نیز ایدهٔ مشترک است.

\textbf{پروژه چه چیز اضافه یا متفاوتی دارد؟} واسط پروژه متنی و قابل‌تعویض است: ادامه‌های موتور، وضعیت قبل و بعد و واژه‌های مربوط به ضعف‌های صفحه را به مدل‌های مختلف می‌دهد. آموزش دامنه‌ای از متن کتاب نیاز اولیه به ساخت مجموعهٔ بزرگ زوج‌های برچسب‌دار را کاهش می‌دهد. ابزار محلی و رابط وب نیز در مخزن پروژه ملموس است؛ این‌ها مزیت مهندسی‌اند، نه اثبات بهتر بودن زبان تولیدی.

\textbf{چه چیز کمتر دارد؟} رابطهٔ مستقیم میان شواهد صفحه و پاسخ هدف در آموزش فعلی وجود ندارد. الگوی کتاب می‌تواند سبک بیان را منتقل کند، اما تضمین نمی‌کند مدل بفهمد همان ایده دقیقاً در کدام موقعیت صدق می‌کند. کنترل محتوای آموخته‌شده و ارزیابی انسانی مقاله نیز معادل مستقیمی در شواهد فعلی پروژه ندارند.

\textbf{آزمایش تعیین‌کننده.} یک مدل پایهٔ ثابت را با سه روش «بدون آموزش»، «آموزش کتاب» و «آموزش زوج‌های موتور\lr{–}توضیح» بسنجید. دادهٔ آزمون را در سطح بازی جدا کنید. تا قبل از این آزمایش، نمی‌توان کاهش نیاز به برچسب را بدون افت کیفیت، دستاورد اثبات‌شده دانست. همچنین کوچک‌تر بودن پروژه نسبت به مدل این مقاله از صرف نام \en{SLM} نتیجه نمی‌شود.
\source{https://arxiv.org/abs/2212.08195}{صفحهٔ آرشیوی و متن مقاله}

\paper{مقایسه با \en{Hebbar et al.}}{Hallucinations on the Board: Tool-Augmented Evaluation of LLM Chess Commentary}{پیش‌چاپ \en{arXiv}، اوت ۲۰۲۶}
\textbf{روش مقاله.} \en{ACT-Eval} متن را به ادعاهای مستقل می‌شکند. ادعاهای قابل‌محاسبه به ابزار صفحه و موتور فرستاده می‌شوند و پوشش ایده‌های تخصصی با مرجع‌های اتمی سنجیده می‌شود. مجموعهٔ ۳۲۵ زوج موقعیت\lr{–}حرکت و مطالعهٔ تطبیق با داوران انسانی ارائه شده است. مبنای مقایسه، چارچوب ارزیابی و بخش اعتبارسنجی انسانی است.

\textbf{وجه مشترک.} هر دو کار فاصلهٔ روانی متن با صحت شطرنجی را جدی می‌گیرند و چند بُعد ارزیابی دارند. اما داشتن چند امتیاز در خروجی داور با بررسی محاسباتی تک‌تک ادعاها یکسان نیست.

\textbf{پروژه چه چیز اضافه یا متفاوتی دارد؟} پروژه یک مولد و برنامهٔ کاربری هم دارد؛ هدف این مقاله عمدتاً ارزیابی است. داور محلی پروژه می‌تواند یک غربال اولیهٔ مستقل از سرویس میزبان باشد و محاسبهٔ امتیاز کل و تطبیق سطرها روشن است. بنابراین گسترهٔ محصول بیشتر است، ولی این به معنی ارزیاب قوی‌تر نیست.

\textbf{چه چیز کمتر دارد؟} تابع ساخت پرامپت داور پروژه فقط پاسخ و مرجع را می‌فرستد. ابزار بررسی صفحه، تفکیک ادعا، پوشش اتمی ایده‌ها، طبقه‌بندی خطای شطرنجی و کالیبراسیون انسانی در این مسیر وجود ندارد. مجموعهٔ منتخب داخلی نیز ۶۰ نمونه دارد. اتکای صرف به متن مرجع می‌تواند هم خطای واقعی را نبیند و هم توضیح درست دیگری را نادرست بشمارد.

\textbf{آزمایش تعیین‌کننده.} خروجی همان مدل‌های پروژه را با داور فعلی، یک ارزیاب ابزارمحور و دو داور شطرنج ارزیابی کنید. موارد اختلاف را دسته‌بندی کنید. این مقاله برای بهبود اعتبار نتیجه‌های پروژه، اولویت بالاتری از افزودن یک مدل مولد دیگر دارد؛ این اولویت قضاوت این سند است، نه نتیجهٔ آزمایش جدید.
\source{https://arxiv.org/abs/2608.04240}{صفحهٔ آرشیوی و متن مقاله}

\paper{مقایسه با \en{Tang et al.}}{Grounded Chess Reasoning in Language Models via Master Distillation}{پیش‌چاپ \en{arXiv}، مارس ۲۰۲۶؛ بازنگری ژوئن}
\textbf{روش مقاله.} راه‌حل موتور به دادهٔ استدلالی زبانی تبدیل می‌شود و مدل کوچک‌تر با آموزش نظارت‌شده و یادگیری تقویتی آموزش می‌بیند. مدل \en{C1} چهار میلیارد پارامتر دارد. تمرکز بر حل مسئلهٔ شطرنج و توضیح راه‌حل است؛ در روش مقاله، جواب نهاییِ حرکت نقش مهمی در پاداش دارد. بخش ۳ و آزمایش‌های حل مسئله مبنای بررسی‌اند.

\textbf{وجه مشترک.} هر دو مسیر می‌خواهند دانش موتور به مدل زبانی نسبتاً کوچک کمک کند. بنابراین «مدل کوچک به همراه موتور و توضیح» به‌تنهایی مرز نوآوری پروژه نیست.

\textbf{پروژه چه چیز اضافه یا متفاوتی دارد؟} هدف پروژه توضیح یک حرکت انجام‌شده در بازی عادی، از جمله اشتباه، است. خروجی کوتاه و مسیر وب برای بازخورد به مبتدی طراحی شده است. مدل ۱٫۷ میلیاردی پروژه از مدل چهار میلیاردی این کار کوچک‌تر است؛ اما هزینهٔ کل شامل اجرای موتور و احتمالاً داور نیز می‌شود و نسبت پارامترها معادل نسبت زمان اجرا نیست.

\textbf{چه چیز کمتر دارد؟} آموزش اختصاصی روی شواهد موتور و توضیح مسئله، بازخورد قابل‌بررسیِ پاسخ نهایی و آزمایش‌های آموزشی روش در پروژه دیده نمی‌شود. آموزش کتاب به تنهایی همان انتقال دانشِ وابسته به موقعیت را انجام نمی‌دهد. در سوی دیگر، درست بودن حرکت نهایی مقاله هم به‌تنهایی صحت همهٔ جمله‌های استدلال را اثبات نمی‌کند.

\textbf{آزمایش تعیین‌کننده.} به‌جای مقایسهٔ مستقیم امتیاز حل مسئله با نمرهٔ داور پروژه، مجموعهٔ مشترکی از حرکت درست و اشتباه بسازید. مدل ۱٫۷ میلیاردی را با توضیح‌های مبتنی بر موتور آموزش دهید و صحت ادعا و پوشش علت اشتباه را بسنجید. این کار مشخص می‌کند کوچک ماندن مدل با دادهٔ هدفمند چه کیفیتی ممکن می‌کند.
\source{https://arxiv.org/abs/2603.20510}{صفحهٔ آرشیوی و نسخهٔ بازنگری‌شده}

\paper{مقایسه با \en{Cui et al.}}{Communicating Chess Strategies in Natural Language}{پیش‌چاپ \en{arXiv}، ژوئیهٔ ۲۰۲۶}
\textbf{روش مقاله.} یک درخت راهبرد جزئی از موتور ساخته و به زبان تبدیل می‌شود. شاخه‌ها واکنش‌های مختلف حریف را لحاظ می‌کنند. ارزیابی بررسی می‌کند دریافت‌کنندهٔ انسانی یا زبانی تا چه اندازه می‌تواند راهبرد را اجرا کند؛ صرف شباهت متن هدف نیست. بخش ۴ و آزمایش‌های انتقال راهبرد مبنا هستند.

\textbf{وجه مشترک.} هر دو، فهم انسانی از پیشنهاد موتور را هدف می‌گیرند. با این حال، توضیح پس از یک حرکت با ارائهٔ سیاست عملی برای چند پاسخ آینده دو مسئلهٔ متفاوت است.

\textbf{پروژه چه چیز اضافه یا متفاوتی دارد؟} توضیح دو یا سه جمله‌ای برای بازخورد فوری پس از خطا دامنه‌ای محدودتر و مناسب یک کاربرد روزمره دارد. مقایسهٔ صریح قبل و بعد، پیامد حرکت انجام‌شده را محور قرار می‌دهد. این محدودسازی می‌تواند بار خواندن را کم کند، اما سرعت و اثر آموزشی آن هنوز اندازه‌گیری نشده است.

\textbf{چه چیز کمتر دارد؟} سه مسیر اصلی موتور معادل درختی نیست که در چند وضعیت آینده برای پاسخ‌های مختلف حریف تصمیم مشخص کند. پروژه ساخت و انتقال راهبرد شاخه‌دار یا آزمون اجرای آن را ندارد. منع نام حرکت و نمادگذاری، برای توضیح کوتاه قابل‌فهم است، اما ممکن است انتقال برنامهٔ دقیق به بازیکن پیشرفته را محدود کند.

\textbf{آزمایش تعیین‌کننده.} دو حالت جدا تعریف کنید: توضیح کوتاه خطا و توضیح ادامهٔ قابل‌اجرا. برای دومی، چند پاسخ مهم حریف و جواب درست به آن‌ها را ارائه دهید و موفقیت کاربر در موقعیت‌های بعدی را بسنجید. ادعای مربی‌بودن پروژه هنگامی قوی‌تر می‌شود که یادگیری یا تصمیم بهتر کاربر اندازه‌گیری شود، نه فقط پسندیده شدن متن.
\source{https://arxiv.org/abs/2607.11486}{صفحهٔ آرشیوی و متن مقاله}

\paper{مقایسه با \en{Zang et al.}}{Automated Chess Commentator Powered by Neural Chess Engine}{\en{ACL 2019}}
\textbf{روش مقاله.} موتور عصبی و مولدهای توضیح به‌صورت مشترک آموزش می‌بینند. بازنمایی صفحه، پیش‌بینی حرکت و ادامهٔ بازی به تولید پنج نوع توضیح کمک می‌کند: توصیف، کیفیت، زمینه، مقایسه و برنامه‌ریزی. مقاله ارزیابی خودکار و انسانی دارد. بخش معرفی معماری و آزمایش‌ها مبنای مقایسه‌اند.

\textbf{وجه مشترک.} هر دو از دانشی فراتر از رشتهٔ خام حرکات استفاده می‌کنند. اختلاف اصلی در محل اتصال موتور و زبان است: بازنمایی داخلی آموخته‌شده در مقاله، در برابر شواهد متنی صریح موتور مستقل در پروژه.

\textbf{پروژه چه چیز اضافه یا متفاوتی دارد؟} برای تولید توضیح لازم نیست یک شبکهٔ بازی‌کننده آموزش داده شود. موتور و مولد را می‌توان مستقل جایگزین کرد و خروجی ویژگی‌ها را بررسی کرد. ابزار دسته‌ای پروژه نیز مقایسهٔ چند مولد را تسهیل می‌کند. این ساختار برای توسعهٔ محصول و آزمایش تعویض مدل انعطاف‌پذیر است.

\textbf{چه چیز کمتر دارد؟} پروژه بازنمایی مشترکی که مستقیماً برای توضیح و پیش‌بینی آموزش دیده باشد ندارد. پوشش پنج نوع توضیح با مدل‌ها و سنجش مستقل هر دسته نیز در شواهد بررسی‌شده دیده نمی‌شود. مسیرهای موتور در پرامپت، جای یادگیری نحوهٔ استفاده از آن‌ها را تضمین نمی‌کنند؛ نمرهٔ پایین مدل اختصاصی همین نیاز را برجسته می‌کند.

\textbf{آزمایش تعیین‌کننده.} آزمون را به انواع توضیح تقسیم کنید و به‌خصوص «علت بد بودن حرکت» را از «بهترین ادامه» جدا بسنجید. همچنین ورودی با و بدون مسیرهای موتور را مقایسه کنید. مدرن‌تر بودن مولد پروژه دلیل برتری بر روش ۲۰۱۹ نیست؛ حتی ادعای سبک‌تر بودن نسبت به کل معماری آن بدون سنجش منابع قابل‌دفاع نیست.
\source{https://aclanthology.org/P19-1597/}{صفحهٔ انتشار و متن کامل مقاله}

\paper{مقایسه با \en{Jhamtani et al.}}{Learning to Generate Move-by-Move Commentary for Chess Games from Large-Scale Social Forum Data}{\en{ACL 2018}}
\textbf{روش مقاله.} بیش از ۲۹۸ هزار زوج حرکت\lr{–}توضیح از حدود ۱۱ هزار بازی گردآوری شده است. مدل از ویژگی‌های حرکت، تهدید و امتیاز برای تولید متن استفاده می‌کند. ارزیابی انسانی گزارش‌شده دربارهٔ شباهت به انسان، محدود به زیرمجموعهٔ توصیف مستقیم حرکت است؛ نباید آن را به تمام توضیح‌های راهبردی تعمیم داد.

\textbf{وجه مشترک.} ویژگی‌های دستی، وضعیت بازی و تولید متن حرکت‌به‌حرکت در هر دو وجود دارند. پس استخراج ویژگی صفحه و افزودن امتیاز موتور به مولد، به‌تنهایی نوآوری تازه‌ای نیست.

\textbf{پروژه چه چیز اضافه یا متفاوتی دارد؟} مسیر فعلی پروژه چند ادامهٔ موتور و قواعد تغییر وضعیت را در یک پرامپت برای مولد عمومی آماده می‌کند. مدل با متن کتاب دامنه‌سازگار می‌شود و تولید کوتاه مبتدی‌محور، سرویس محلی و رابط وب دارد. مزیت احتمالی، شروع از دانش زبانی ازپیش‌آموخته و انعطاف پیاده‌سازی است؛ این مزیت در آزمون مشترک هنوز سنجیده نشده است.

\textbf{چه چیز کمتر دارد؟} در مقایسه با مجموعهٔ بزرگ جفت‌شدهٔ مقاله، کتاب‌های بدون اتصال به وضعیت، نظارت مستقیم بسیار کمتری برای وظیفهٔ هدف فراهم می‌کنند. مجموعهٔ ۶۰تایی پروژه، دادهٔ آزمون است و نباید با ۲۹۸ هزار نمونهٔ آموزش مقایسهٔ عددی مستقیم شود. داوری انسانی و خط مبناهای کلاسیک مقاله نیز در پروتکل فعلی پروژه بازتولید نشده‌اند.

\textbf{آزمایش تعیین‌کننده.} یک زیرمجموعهٔ تمیز از دادهٔ جفت‌شده را به آموزش بیفزایید و بازی‌های آزمون را جدا نگه دارید. مقایسه با یک روش قالبی و بازیابی توضیح مشابه، نشان می‌دهد مدل کوچک واقعاً چه ارزشی فراتر از روش ساده می‌افزاید. کیفیت توصیف ظاهری را جدا از صحت علت تاکتیکی گزارش کنید.
\source{https://aclanthology.org/P18-1154/}{صفحهٔ انتشار و متن کامل مقاله}

\paper{مقایسه با \en{dataChess / CARLSy}}{Meet dataChess and CARLSy: Towards the Explanation of Chess Plays}{نسخهٔ ناشناس در \en{OpenReview}؛ سال و پذیرش قطعی در نسخهٔ بررسی‌شده تأیید نشد}
\textbf{روش منبع.} دادهٔ بازی و توضیح پاک‌سازی می‌شود و اطلاعاتی مانند جای مهره‌ها، حمله‌ها و دستهٔ طول توضیح به آن افزوده می‌شود. مدل برای توضیح کیفیت حرکت آموزش می‌بیند و ارزیابی خودکار و انسانی گزارش می‌شود. ساخت داده در بخش ۳ و مدل و ارزیابی در ادامهٔ متن مبنا هستند.

\textbf{وجه مشترک.} هر دو پروژه اطلاعات ساختاری صفحه را برای کمک به مدل زبانی آماده می‌کنند و تولید توضیح آموزشی را هدف دارند. این منبع برای طراحی دادهٔ آموزشی از نظر عملی نزدیک است.

\textbf{پروژه چه چیز اضافه یا متفاوتی دارد؟} در کد \en{BlunderGuard} تحلیل قبل و بعد، سه ادامهٔ برتر موتور، وضعیت مات و ویژگی‌های دو طرف در یک رکورد جمع می‌شوند. بنابراین شواهد صریح جست‌وجوی موتور و اثر حرکت در واسط پروژه برجسته است. متن بررسی‌شدهٔ \en{CARLSy} چنین مسیر دقیقِ هم‌ارزی را نشان نمی‌دهد؛ این گزاره به معنی نبود هرگونه موتور در همه نسخه‌های آن نیست.

\textbf{چه چیز کمتر دارد؟} دادهٔ جفت‌شدهٔ پاک‌سازی‌شده و آموزش مستقیم بر زمینهٔ صفحه از نقاط قوت منبع است. پروژه مجموعهٔ بزرگی با همین ساختار آموزشی و ارزیابی انسانی متناظر نشان نمی‌دهد. متن کتاب می‌تواند قواعد عمومی را یاد بدهد، اما برای تشخیص اینکه توضیح به همین حرکت مربوط است، دادهٔ جفت‌شده مسیر مستقیم‌تری دارد.

\textbf{آزمایش تعیین‌کننده.} قالب آموزشی شامل صفحه، حرکت، ویژگی‌ها، مسیر موتور و توضیح تأییدشده بسازید. سپس اثر هر نوع زمینه را جدا حذف و اندازه‌گیری کنید. در متن پایان‌نامه این منبع باید «نسخهٔ ناشناس» معرفی شود؛ نام نویسنده، سال پذیرش یا محل انتشار را از روی حدس کامل نکنید.
\source{https://openreview.net/pdf/f3eee18d4e2efc80301095c1ead34005ef8813c0.pdf}{نسخهٔ کامل ناشناس در آرشیو}

\paper{مقایسه با مقالهٔ فنی \en{Tredence}}{Fine-tuning a chess commentary model}{وبلاگ فنی؛ تاریخ انتشار در متن بازیابی‌شده مشخص نیست}
\textbf{روش منبع.} مقاله تنظیم مدل \en{Llama 3.3 70B Instruct} برای توضیح شطرنج را با دادهٔ \en{ChessGPT} و قالب مکالمه در \en{Databricks} شرح می‌دهد. نمونه‌هایی از قبل و بعد و مقایسه با چند مدل ارائه می‌شود. نتیجهٔ کیفی نویسنده بر انسانی‌تر شدن لحن تأکید دارد؛ این را نباید شاهد اثبات صحت شطرنجی دانست.

\textbf{وجه مشترک.} هر دو از آموزش تکمیلی مدل زبانی برای توضیح شطرنج استفاده می‌کنند. تفاوت نوع دادهٔ آموزشی مهم‌تر از صرف عبارت «فاین‌تیون» است: متن کتاب پروژه با نمونهٔ مکالمهٔ وابسته به بازی یکسان نیست.

\textbf{پروژه چه چیز اضافه یا متفاوتی دارد؟} مولد ۱٫۷ میلیاردی در برابر مدل ۷۰ میلیاردی این مقاله قرار می‌گیرد و مسیر \en{GGUF} محلی در کد دارد. پروژه همچنین مجموعه‌آزمون ثبت‌شده و پنج بُعد امتیازدهی دارد. از نظر گزارش کمیِ قابل‌ردیابی، این شواهد از نتیجه‌گیری صرفاً مبتنی بر لحنِ نمونه‌های وبلاگ ساخت‌یافته‌تر است.

\textbf{چه چیز کمتر دارد؟} مسیر آموزش فعلی پروژه دستور و پاسخ هدفِ وابسته به بازی را مستقیماً آموزش نمی‌دهد. مزیت آموزشی قالب مکالمه باید برای مدل کوچک پروژه آزمایش شود. وجود مدل بزرگ در وبلاگ نیز دلیل کافی برای نتیجه‌گیری دربارهٔ کیفیت بهتر آن در وظیفهٔ شما نیست؛ آزمون مشترک وجود ندارد.

\textbf{آزمایش تعیین‌کننده.} سبک، صحت و کاربرد آموزشی را جدا بسنجید. پاسخ «احساسی‌تر» یا «شبیه انسان» ممکن است غلط باشد. با همان مدل کوچک، آموزش کتاب و آموزش مکالمهٔ مستند به موتور را مقایسه کنید. این منبع برای الهام پیاده‌سازی مفید است، اما نباید در ادعای برتری علمی هم‌وزن مقالهٔ دارای مطالعهٔ انسانی معتبر قرار گیرد.
\source{https://www.tredence.com/blog/finetuning-a-chess-commentary-model}{متن مقالهٔ فنی و پیوندهای نمونه‌ها}

\clearpage
\section{محدودیت‌های کد که بر مقایسه اثر می‌گذارند}
این موارد برای دقیق بودن توصیف قابلیت فعلی ثبت شده‌اند؛ در این کار تغییری در آن‌ها اعمال نشده است.
\begin{enumerate}
\item \textbf{مشخصات مدل یکدست نیست.} \en{README} مدل ۱٫۶ میلیاردی و \en{Q4\_K\_M} را ذکر می‌کند، ولی نوت‌بوک پایهٔ ۱٫۷ میلیاردی و تبدیل \en{Q8\_0} دارد. شناسهٔ نتیجهٔ اختصاصی هم \en{chess-q8\_0} است. سند حاضر به نوت‌بوک و خروجی منتخب اولویت می‌دهد؛ اتصال دقیق نتیجه به وزن مدل نیازمند هش و شناسهٔ ثبت‌شده است.
\item \textbf{ادعای \en{QLoRA} نیاز به اصلاح یا شاهد دارد.} نوت‌بوک \en{LoRA} را نشان می‌دهد، اما در بارگذاری آموزشیِ بررسی‌شده تنظیم بارگذاری چهار بیتی دیده نمی‌شود. کمّی‌سازی بعد از آموزش به \en{GGUF} به‌تنهایی \en{QLoRA} نیست.
\item \textbf{مسیر مستقیم خطای قالب دارد.} قالب \en{SINGLE\_METHOD} فیلدهای اضافی مانند \en{missed\_opportunity} می‌خواهد، ولی \en{single\_method} همه را نمی‌فرستد. فراخوانی مستقل با مولد ساختگی پیش از هر تماس مدل، خطای \en{KeyError} برای این فیلد داد. بنابراین نباید صرف وجود مسیر \en{/single} را شاهد اجرای موفق فعلی آن شمرد؛ مسیر تولید دسته‌ای سازندهٔ پرامپت جدا دارد.
\item \textbf{قواعد شطرنجی کامل نیستند.} تشخیص مهرهٔ بی‌دفاع بر وجود مهاجم و نبود مدافع تکیه دارد و تحلیل کامل تعویض نیست. در محاسبهٔ حمایت پیاده، جهت مقایسه برای دو رنگ یکسان است. در طبقه‌بندی مات، مقدار بولی با نام رنگ مقایسه می‌شود؛ از نظر نوع داده، جهت درست پیامد مات را بیان نمی‌کند. این دو مورد آخر از خواندن کد است و با آزمون جامع شطرنجی بررسی نشده‌اند.
\item \textbf{سرعت و سبک بودن کل سامانه اندازه‌گیری نشده است.} تابع موتور پارامتر زمان می‌گیرد، اما از محدودیت عمق ۱۷ استفاده می‌کند؛ تنظیم فعلی ۱۴ رشته و ۲۰۴۸ مگابایت هش است. اجرای محلی ممکن است، ولی «بلادرنگ روی سخت‌افزار ضعیف» از این کد نتیجه نمی‌شود.
\end{enumerate}
\pathref{README.md | notebooks/chess.ipynb | app/chat.py | app/prompt.py}
\pathref{app/chess/stockfish_analyzer.py | app/chess/position_analysis.py}

\clearpage
\section{چه ادعایی دربارهٔ پروژه قابل دفاع است؟}
\subsection*{نقاط قوت فعلی}
\begin{itemize}
\item مسیر مولد محلی کمّی‌شده و امکان استفاده از موتور مستقل؛ مزیت استقرار و کنترل وابستگی‌ها.
\item واسط صریح شامل قبل و بعد، مسیرهای موتور و ویژگی‌های صفحه؛ امکان بازرسی و آزمایش حذف اجزا.
\item تمرکز روشن بر توضیح کوتاه علت اثر حرکت برای مبتدی؛ هدف محصولی مشخص.
\item زیرساخت مقایسهٔ چند مدل، تطبیق سطرها و جداسازی امتیاز روانی از توافق واقعی با مرجع.
\end{itemize}
\subsection*{ادعاهایی که شواهد کافی ندارند}
دقت بالاتر از مقاله‌ها، کاهش قطعی توهم، اثربخشی آموزشی، برتری روش زنجیره‌ای، بی‌اثر بودن کمّی‌سازی بر کیفیت، برتری آموزش کتاب نسبت به مدل پایه و عملکرد بلادرنگ روی دستگاه ضعیف هنوز اثبات نشده‌اند. همچنین فارسی بودن این سند به معنی وجود سامانهٔ مربی فارسیِ ارزیابی‌شده نیست؛ پرامپت‌های بررسی‌شده انگلیسی‌اند.

\subsection*{صورت‌بندی پیشنهادی برای گزارش یا دفاع}
«\en{BlunderGuard} یک سامانهٔ قابل‌استقرار محلی برای توضیح کوتاه حرکت شطرنج است که مدل زبانی کوچکِ سازگارشده با متون شطرنج را با مسیرهای موتور و ویژگی‌های صریح تغییر صفحه ترکیب می‌کند. ارزیابی داخلی، روانی متن را در کنار ضعف توافق با مرجع نشان می‌دهد؛ بنابراین سهم فعلی، طراحی و بررسی تجربی یک مسیر کم‌حجم برای توضیح مبتدی‌محور است و برتری صحت بر روش‌های پیشین ادعا نمی‌شود.»

\subsection*{فرصت پژوهشی بعدی؛ فرضیه، نه دستاورد فعلی}
پرسش قوی‌تر این است: «آیا یک مدل حداکثر دو میلیارد پارامتری با آموزش زوج‌های مبتنی بر موتور و بررسی ادعا، می‌تواند روی سخت‌افزار مشخص، توضیح دقیق و کوتاه را با هزینهٔ کمتر ارائه کند؟» پاسخ مثبت فقط پس از مقایسهٔ مشترک و اندازه‌گیری منابع، سهم علمی قابل‌دفاع می‌سازد. صرف اتصال موتور به مدل یا استفاده از کتاب، برای ادعای نخستین‌بودن کافی نیست.

\clearpage
\section{برنامهٔ پیشنهادی برای بستن فاصله‌ها}
\begin{enumerate}
\item \textbf{اول: مبنای آزمایش را قابل‌اعتماد کنید.} خطای قالب، جهت مات و قواعد وابسته به رنگ را بررسی و اصلاح کنید. شناسه و هش مدل، نسخهٔ موتور، پرامپت و تنظیم اجرا را ذخیره کنید. تا این مرحله روشن نشود، اختلاف امتیازها ممکن است از اتصال اجزا باشد.
\item \textbf{دوم: اثر آموزش را جدا کنید.} مدل پایه و مدل آموزش‌دیده را با ورودی، قالب مکالمه، طول خروجی و کمّی‌سازی یکسان مقایسه کنید. سپس آموزش کتاب را در برابر آموزش جفت‌شده و ترکیب این دو بسنجید. جداسازی داده در سطح بازی و حذف تکرار لازم است.
\item \textbf{سوم: اثر شواهد را اندازه بگیرید.} حالت‌های بدون موتور، فقط امتیاز، امتیاز و مسیرها، و مجموعهٔ کامل ویژگی‌ها را مقایسه کنید. تعداد مسیرها و روش مستقیم\lr{/}زنجیره‌ای نیز باید با بودجهٔ محاسباتی مشخص مقایسه شوند.
\item \textbf{چهارم: ارزیابی را به واقعیت صفحه وصل کنید.} با الهام از \en{ACT-Eval}، ادعاهای سادهٔ صفحه و تاکتیک را به ابزار بسپارید و پوشش ایدهٔ اصلی را جدا نمره دهید. بخشی از خروجی‌ها را دو داور شطرنج مستقل بررسی کنند و اختلاف داورها گزارش شود.
\item \textbf{پنجم: دادهٔ آزمون را متنوع کنید.} گروه‌بندی سطح بازیکن را با برچسب نوع خطا، تاکتیک، مرحلهٔ بازی و دشواری توضیح تکمیل کنید. نتایج را با بازهٔ اطمینان و نمونه‌های شکست منتشر کنید. از ارزیابی فقط روی نمونه‌های انتخاب‌شدهٔ موفق پرهیز کنید.
\item \textbf{ششم: مزیت عملی را اندازه بگیرید.} حافظهٔ اوج، اندازهٔ فایل، زمان موتور و زمان تولید را جدا ثبت کنید؛ برای تأخیر، میانه و صدک ۹۵ مناسب‌اند. آزمون یادگیری کاربر باید نشان دهد توضیح کوتاه واقعاً به تصمیم بهتر کمک می‌کند.
\end{enumerate}

\textbf{اولویت خواندن:} برای معماری، \en{Lee} و \en{Kim}؛ برای آموزش، \en{Master Distillation} و \en{CARLSy}؛ برای اعتبار ارزیابی، \en{ACT-Eval}؛ برای تبدیل توضیح به آموزش عملی، \en{Cui}. این ترتیب بر اساس نیاز فعلی پروژه پیشنهاد شده است.

\clearpage
\section{فهرست منابع و شواهد قابل پیگیری}
پیوند عنوان هر منبع به صفحهٔ اصلی یا نسخهٔ کامل آن متصل است. تاریخ دسترسی همهٔ منابع: \en{2026-09-13}. وضعیت پیش‌چاپ به معنی تأیید داوری علمی نیست.
\begin{enumerate}
\item \en{Kim, J. et al. (2025).} \href{https://aclanthology.org/2025.naacl-long.481/}{\en{Concept-guided Chess Commentary (CCC).}} \en{NAACL}. روش: بخش ۳؛ ارزیابی و محدودیت‌ها: بخش‌های ۴ و ۷.
\item \en{Lee, A.; Wu, D.; Dinan, E.; Lewis, M. (2022).} \href{https://arxiv.org/abs/2212.08195}{\en{Improving Chess Commentaries by Combining Language Models with Symbolic Reasoning Engines.}} \en{arXiv:2212.08195}.
\item \en{Hebbar, S. A. et al. (2026).} \href{https://arxiv.org/abs/2608.04240}{\en{Hallucinations on the Board.}} \en{arXiv:2608.04240}. چارچوب \en{ACT-Eval} و اعتبارسنجی انسانی.
\item \en{Tang, Z. et al. (2026).} \href{https://arxiv.org/abs/2603.20510}{\en{Grounded Chess Reasoning in Language Models via Master Distillation.}} \en{arXiv:2603.20510v2}.
\item \en{Cui, L.; Ling, C. K.; Ng, H. T. (2026).} \href{https://arxiv.org/abs/2607.11486}{\en{Communicating Chess Strategies in Natural Language.}} \en{arXiv:2607.11486}.
\item \en{Zang, H.; Yu, Z.; Wan, X. (2019).} \href{https://aclanthology.org/P19-1597/}{\en{Automated Chess Commentator Powered by Neural Chess Engine.}} \en{ACL}, صفحات ۵۹۵۲ تا ۵۹۶۱.
\item \en{Jhamtani, H. et al. (2018).} \href{https://aclanthology.org/P18-1154/}{\en{Learning to Generate Move-by-Move Commentary for Chess Games from Large-Scale Social Forum Data.}} \en{ACL}, صفحات ۱۶۶۱ تا ۱۶۷۱.
\item نویسندگان ناشناس. \href{https://openreview.net/pdf/f3eee18d4e2efc80301095c1ead34005ef8813c0.pdf}{\en{Meet dataChess and CARLSy: Towards the Explanation of Chess Plays.}} نسخهٔ پنج‌صفحه‌ای بررسی‌شده در \en{OpenReview}.
\item \en{Tredence.} \href{https://www.tredence.com/blog/finetuning-a-chess-commentary-model}{\en{Fine-tuning a chess commentary model.}} مقالهٔ فنی، بخش‌های تنظیم آزمایش و نتیجه‌گیری.
\end{enumerate}
\textbf{شواهد مخزن:} مسیرها در بخش‌های ۲، ۳ و ۱۳ ثبت شده‌اند. آموزش از سلول‌های کد نوت‌بوک و اعداد از فایل خلاصه بازخوانی شدند. متن گزارش برای ادعاهای آموزشی و منشأ داده استفاده شده، اما در تعارض با کد، به‌عنوان شاهد قطعی کارکرد پذیرفته نشده است. فایل منبع همین سند کنار \en{PDF} نگهداری می‌شود تا اصلاح و به‌روزرسانی آن ممکن باشد.
\end{document}
